\documentclass[11pt]{article}
\usepackage[a4paper,margin=1in]{geometry}
\usepackage[T1]{fontenc}
\usepackage{lmodern}
\usepackage{microtype}
\usepackage{amsmath}
\usepackage{booktabs}
\usepackage{fancyhdr}
\usepackage[hidelinks]{hyperref}

\pagestyle{fancy}
\fancyhf{}
\renewcommand{\headrulewidth}{0pt}
\fancyfoot[L]{ET23BTIT133}
\fancyfoot[R]{\thepage}

\begin{document}
\begin{center}
{\LARGE\bfseries Supervised Machine Learning Approach for\\
WSD in Spoken News\par}
\vspace{0.8em}
{\Large Architecture and Methodology\par}
\vspace{0.8em}
{\large Sree Saicharan Vadapalli\\ ET23BTIT133\par}
\end{center}
\vspace{1em}

\section{Overview}

News transcripts are full of words like \emph{bank}, \emph{interest},
\emph{plant} and \emph{run} whose meaning depends on the sentence around them.
This project treats word sense disambiguation (WSD) as a supervised
classification problem. For every target word we train a separate classifier
(a ``word expert'') that maps the context of an occurrence to a WordNet synset.
Two classifiers are compared, a linear SVM and multinomial Naive Bayes, each
with two context window sizes, $W = \pm 2$ and $W = \pm 5$.

The pipeline has three stages, one per source file:

\begin{center}
\begin{tabular}{ll}
\toprule
File & Responsibility \\
\midrule
\texttt{src/preprocessing.py} & SemCor extraction, POS tags, lemmas,
stopwords \\
\texttt{src/models.py} & Feature extraction and the two classifiers \\
\texttt{src/pipeline.py} & Cross-validated training and testing, metrics \\
\bottomrule
\end{tabular}
\end{center}

\section{Dataset}

The dataset is built from SemCor \cite{semcor}, read through NLTK. SemCor was
picked over NLTK's Senseval corpus for two reasons: its labels are WordNet
synsets, which is what we want to predict, and it covers the words named in
the problem. NLTK's Senseval data only has \emph{hard}, \emph{interest},
\emph{line} and \emph{serve}, labelled with its own sense inventory.

Besides the four words from the problem statement we added \emph{state},
\emph{case}, \emph{rate} and \emph{power}. They are common in news, have
several senses, and have enough tagged examples in SemCor to train on.

Extraction walks every sense-tagged sentence and keeps each chunk whose WordNet
lemma is one of the targets. For each hit we store the target word, its synset,
its token position and the full sentence, then POS-tag the sentence with NLTK's
perceptron tagger. This gives 1,177 labelled occurrences, saved to
\texttt{data/annotated/semcor\_wsd.csv}. Senses with fewer than 5 examples are
dropped before training since they cannot appear in every fold of a 5-fold
split, leaving 1,092 occurrences.

\begin{center}
\begin{tabular}{lrrr}
\toprule
Word & Occurrences & Senses & Majority sense share \\
\midrule
bank & 34 & 2 & 0.59 \\
case & 133 & 5 & 0.48 \\
interest & 140 & 7 & 0.41 \\
plant & 107 & 3 & 0.59 \\
power & 99 & 6 & 0.57 \\
rate & 116 & 3 & 0.59 \\
run & 185 & 12 & 0.36 \\
state & 278 & 6 & 0.39 \\
\bottomrule
\end{tabular}
\end{center}

One caveat: SemCor is written text from the Brown corpus, not speech. Raw
transcripts have no punctuation, so the punctuation tags that show up in the POS
window here would be missing, and the POS tagger itself was trained on
punctuated text and will be less accurate on them.

\section{Feature vectors}

Take a sentence $w_1, \dots, w_n$ with POS tags $p_1, \dots, p_n$ and a target
word at position $t$. For a window size $W$ the feature vector has three parts:

\begin{enumerate}
\item \textbf{Target POS}: the tag $p_t$, one-hot encoded. This alone
separates noun and verb senses such as \emph{plant.n.01} and
\emph{plant.v.01}.
\item \textbf{Neighbouring POS}: the tags $p_{t+k}$ for
$k \in \{-W, \dots, -1, 1, \dots, W\}$, each as its own positional feature,
so \texttt{pos\_-1=JJ} and \texttt{pos\_+1=JJ} are different features.
\item \textbf{Window bag-of-words}: counts of the lemmas of $w_{t+k}$ over the
same offsets. Tokens are lowercased, stopwords and non-alphabetic tokens are
removed, and the rest are lemmatized with the WordNet lemmatizer using the
token's POS tag.
\end{enumerate}

Putting it together:
\[
\mathbf{x}_t = \mathrm{onehot}(p_t) \;\oplus
\bigoplus_{0<|k|\le W} \mathrm{onehot}_k(p_{t+k}) \;\oplus
\sum_{0<|k|\le W} \mathrm{onehot}\big(\mathrm{lemma}(w_{t+k})\big)
\]
where $\oplus$ is concatenation.

The window is counted on raw token positions, so $W$ means the same thing for
the POS and bag-of-words parts. A side effect is that at $W = \pm 2$ the
bag-of-words often holds only one or two words, since stopwords take up some of
the slots.

A real example from the dataset, with \emph{interest} as the target:

\begin{quote}
\small
\ldots considering the widespread \textbf{interest} in the election , the
number of voters \ldots
\end{quote}

\noindent At $W = \pm 2$ this becomes

{\small
\begin{verbatim}
{pos_0=NN, pos_-2=DT, pos_-1=JJ, pos_+1=IN, pos_+2=DT, bow_widespread=1}
\end{verbatim}
}

\noindent At $W = \pm 5$ it also picks up
\texttt{bow\_consider}, \texttt{bow\_election} and the tags out to
\texttt{pos\_-5} and \texttt{pos\_+5}.

The feature dicts are turned into vectors with scikit-learn's
\texttt{DictVectorizer}. It sits inside the model pipeline, so the vocabulary is
learned from the training folds only.

\section{Classifiers}

Both classifiers are scikit-learn defaults \cite{sklearn}; we did no
hyperparameter tuning.

\paragraph{Multinomial Naive Bayes} picks the sense $s$ that maximizes
\[
\hat{s} = \arg\max_s \; \log P(s) + \sum_j x_j \log P(f_j \mid s),
\]
where $P(f_j \mid s)$ is estimated from feature counts with Laplace smoothing
($\alpha = 1$). It assumes features are independent given the sense, which is
clearly not true for overlapping POS and word features, but it is fast and works
well on small data.

\paragraph{Linear SVM} (\texttt{SVC} with a linear kernel, $C = 1$) learns a
separating hyperplane $f(\mathbf{x}) = \mathbf{w}^\top \mathbf{x} + b$ by
minimizing
\[
\frac{1}{2}\lVert \mathbf{w} \rVert^2 + C \sum_i
\max\big(0,\, 1 - y_i f(\mathbf{x}_i)\big).
\]
Multiple senses are handled one-vs-one, which is what \texttt{SVC} does by
default. A linear kernel is the usual choice for sparse, high-dimensional text
features.

\section{Evaluation protocol}

With 34 to 278 examples per word, a single train/test split would leave some
test sets with fewer than 10 examples. Instead we use stratified 5-fold
cross-validation (shuffled, fixed seed) through \texttt{cross\_val\_predict}.
Every occurrence is predicted exactly once, by a model that never saw it in
training, and the metrics are computed on those predictions.

We report accuracy and macro-averaged precision, recall and F1:
\[
\mathrm{P}_s = \frac{TP_s}{TP_s + FP_s}, \quad
\mathrm{R}_s = \frac{TP_s}{TP_s + FN_s}, \quad
\mathrm{F1}_s =
\frac{2\,\mathrm{P}_s \mathrm{R}_s}{\mathrm{P}_s + \mathrm{R}_s},
\quad
\text{Macro F1} = \frac{1}{|S|} \sum_{s \in S} \mathrm{F1}_s.
\]
Sense distributions are skewed, so accuracy alone rewards a model that just
predicts the most common sense. Macro F1 weighs every sense equally and shows
whether rare senses are actually being learned.

\section{Running it}

{\small
\begin{verbatim}
pip install -r requirements.txt
python -m nltk.downloader semcor wordnet stopwords averaged_perceptron_tagger_eng
python src/preprocessing.py
python src/pipeline.py
\end{verbatim}
}

\noindent The first script writes the dataset CSV. The second prints a summary
table and writes per-word metrics and per-occurrence predictions to
\texttt{results/}. Results are discussed in the Evaluation Report.

\begin{thebibliography}{9}
\bibitem{semcor} G. A. Miller, C. Leacock, R. Tengi and R. T. Bunker. A
semantic concordance. \emph{Proceedings of the ARPA Human Language Technology
Workshop}, 1993.
\bibitem{sklearn} F. Pedregosa et al. Scikit-learn: Machine learning in
Python. \emph{Journal of Machine Learning Research}, 12:2825--2830, 2011.
\end{thebibliography}

\end{document}
